\documentclass[10pt,twocolumn]{article}

% Packages
\usepackage[a4paper,margin=0.7in]{geometry}
\usepackage{graphicx}
\usepackage{booktabs}
\usepackage{array}
\usepackage{caption}
\usepackage{hyperref}
\usepackage{titlesec}
\usepackage{parskip}
\usepackage{amssymb}

\setlength{\columnsep}{0.28in}
\setlength{\parskip}{3pt}

\graphicspath{{figures/}}

\hypersetup{
    colorlinks=true,
    urlcolor=blue,
    linkcolor=black
}

\captionsetup{
    font=small,
    labelfont=bf
}

\titlespacing*{\section}{0pt}{6pt}{2pt}
\titleformat{\section}
{\normalfont\large\bfseries}
{\thesection.}{0.5em}{}

\begin{document}

\twocolumn[{%
\centering
{\LARGE\bfseries MIT 805 Big Data Semester Project -- Part 1}\par
\vspace{3pt}

{\large NYC TLC High-Volume For-Hire Vehicle Trip Records}\par
\vspace{6pt}

{\normalsize
\textbf{Group [NN]} \quad | \quad
[Member 1 Name] ([Student No.]) \quad
[Member 2 Name] ([Student No.])
}\par
\vspace{3pt}

{\small
GitHub:
\url{https://github.com/[user]/[repo]}
}\par

\vspace{6pt}
\rule{\textwidth}{0.4pt}
\vspace{8pt}
}]

% ============================================================
\section{Dataset and source}

This project uses the New York City Taxi and Limousine Commission
(NYC TLC) High-Volume For-Hire Vehicle (HVFHV) Trip Record Data.
HVFHV records contain trips reported to the TLC by licensed
high-volume for-hire vehicle bases. The dataset was selected because
it provides a large public source of urban mobility data and includes
information that can be used to study when and where trips occur,
trip characteristics, fares and operator activity.

The data is published by the NYC TLC as monthly Parquet files.
Fields include request, pickup and drop-off times, pickup and
drop-off location IDs, trip distance, trip time, base passenger
fare, driver pay, tips, tolls, operator codes and service flags
\cite{tlc,dict}.

The data files are not stored in the GitHub repository because of
their size. Instead, the repository provides the source and the
notebook used to identify and access the required files.

% ============================================================
\section{Dataset size}

The sizes of the available HVFHV files were checked programmatically
using the file information returned by the TLC server. This was done
instead of estimating the size of the dataset.

The raw collection contains \textbf{89 monthly files} from
February 2019 to June 2026, with a combined compressed size of
\textbf{37.58 GB}. From this collection, a \textbf{12.00 GB}
working dataset consisting of 26 files was selected. A smaller
\textbf{3.36 GB} processing dataset consisting of 7 files was then
defined for the distributed PySpark work in Part 2.

These sizes satisfy the project requirements of 25--40 GB for the
raw dataset, at least 12 GB for the working dataset and at least
3 GB for the processing dataset. The three levels are shown in
Table~\ref{tab:scale}.

For Part 1, exploratory analysis was carried out in pandas using the
June 2026 file. This monthly sample contains
\textbf{20,775,868 rows and 25 columns}. The one-month EDA sample
should not be confused with the 3.36 GB processing dataset, which
will be used for the distributed analysis in Part 2.

% ============================================================
\section{Data quality and preparation}

The original June 2026 data was kept unchanged and a separate copy
was created for cleaning. Basic checks were first performed for
missing values, duplicates and unusual numerical values.

Several issues were found. There were \textbf{4,844 trips with zero
distance}, \textbf{6,750 trips with a negative base fare} and
\textbf{2 trips with a non-positive duration}. Some very large
distance values were also present. No exact duplicate rows were
found.

Missing values were not equally distributed across the columns. For
example, \texttt{originating\_base\_num} was missing for 26.45\% of
the sample. This is important because missingness in individual
fields can limit the types of comparisons that can be made.

A set of simple plausibility rules was used for the analytical copy.
Trips were retained where duration was greater than 0 and no more
than 1,440 minutes, distance was greater than 0 and no more than
100 miles, and base fare was between \$0 and \$500.

After applying these rules, \textbf{14,421 records (0.07\%)} were
excluded. The cleaned dataset therefore contained
\textbf{20,761,447 trips}.

There are also limitations that cleaning cannot solve. The data
represents HVFHV activity rather than all transport in New York
City, and the June 2026 EDA only represents one month. The results
should therefore not be generalised to all NYC transport or to
seasonal patterns without further analysis.

% ============================================================
\section{Exploratory data analysis}

The cleaned June 2026 sample had an average trip distance of
\textbf{4.98 miles} and an average trip duration of
\textbf{20.36 minutes}. The median values were lower at 2.91 miles
and 16.18 minutes respectively. The average base passenger fare was
\textbf{\$28.75}, compared with a median of \$20.13.

Trip distance and base fare had a strong positive correlation. The
correlation was approximately \textbf{0.86} before cleaning and
remained approximately \textbf{0.86} afterwards. This suggests that
the relatively small number of excluded observations had little
effect on the overall distance--fare relationship.

Trip activity also varied by time. Saturday recorded
\textbf{3,215,911 trips}, while Wednesday recorded
\textbf{2,695,792}. Trip characteristics changed across the day as
well. For example, trips starting around 05:00 averaged approximately
7.13 miles, while trips at 15:00 averaged approximately 4.70 miles
and 23.97 minutes. These differences show that trip demand and trip
characteristics are not constant throughout the day or week.

The spatial analysis also showed differences between locations.
LaGuardia Airport was the largest pickup zone in the sample with
\textbf{421,225 pickups}, followed by JFK Airport with
\textbf{356,638}. At borough level, approximately 7.14 million
pickups occurred in Manhattan, 5.72 million in Brooklyn, 4.77
million in Queens, 2.79 million in the Bronx and 0.35 million in
Staten Island.

The operator field also allowed the sample to be compared by
company. Uber accounted for approximately \textbf{73.4\%} of the
June 2026 trips and Lyft for approximately \textbf{26.6\%}.
Figures showing the main temporal and spatial results are included
in the Appendix.

% ============================================================
\section{The 7 Vs of Big Data}

\textit{[Partner section to be inserted here. The final version will
discuss Volume, Velocity, Variety, Veracity, Variability,
Visualization and Value using evidence from the HVFHV dataset.]}

% ============================================================
\section{Value and limitations}

The HVFHV data can be useful for both operational and planning
purposes. Hourly and daily trip patterns can help identify periods
of high demand, while pickup-zone analysis can show where demand is
concentrated. This could support decisions around driver or vehicle
availability, airport demand and transport planning.

The fare and driver-pay fields also make it possible to study the
financial characteristics of trips, while operator codes allow
comparisons between the major HVFHV services. Over a longer time
period, the dataset could also be used to examine how mobility
patterns change.

There are, however, important limitations. The dataset does not
represent all forms of transport in New York City. The records are
also de-identified, which prevents analysis of individual riders or
drivers over time. Location is provided through taxi-zone IDs rather
than exact GPS coordinates, and some fields contain missing or
unusual values. These limitations need to be considered when
interpreting the results.

% ============================================================
\begin{thebibliography}{9}
\small

\bibitem{tlc}
NYC Taxi \& Limousine Commission.
\emph{TLC Trip Record Data}.
\url{https://www.nyc.gov/site/tlc/about/tlc-trip-record-data.page}

\bibitem{dict}
NYC Taxi \& Limousine Commission.
\emph{High Volume FHV Trip Records Data Dictionary}.
\url{https://www.nyc.gov/assets/tlc/downloads/pdf/data_dictionary_trip_records_hvfhs.pdf}

\bibitem{opendata}
NYC Open Data.
\emph{Frequently Asked Questions}.
\url{https://opendata.cityofnewyork.us/faq/}

\end{thebibliography}

% ============================================================
\clearpage
\appendix

\section*{Appendix}

\setcounter{table}{0}
\renewcommand{\thetable}{A\arabic{table}}

\setcounter{figure}{0}
\renewcommand{\thefigure}{A\arabic{figure}}

% ============================================================
\begin{table*}[t]
\centering
\caption{Dataset sizes used in the project.}
\label{tab:scale}

\begin{tabular}{lrrl}
\toprule
\textbf{Dataset} &
\textbf{Files} &
\textbf{Size (GB)} &
\textbf{Requirement} \\
\midrule

Raw &
89 &
37.58 &
25--40 GB \checkmark \\

Working &
26 &
12.00 &
$\geq$ 12 GB \checkmark \\

Processing &
7 &
3.36 &
$\geq$ 3 GB \checkmark \\

\bottomrule
\end{tabular}
\end{table*}

% ============================================================
\begin{table*}[t]
\centering
\caption{Examples of fields available in the HVFHV dataset.}
\label{tab:dict}
\small

\begin{tabular}{
>{\raggedright\arraybackslash}p{4.2cm}
>{\raggedright\arraybackslash}p{2.0cm}
>{\raggedright\arraybackslash}p{9.6cm}
}

\toprule
\textbf{Field} &
\textbf{Type} &
\textbf{Description} \\
\midrule

\texttt{hvfhs\_license\_num} &
Categorical &
High-volume for-hire service licence code \\

Request / pickup / drop-off times &
Temporal &
Timestamps describing stages of a trip \\

\texttt{PULocationID}, \texttt{DOLocationID} &
Spatial &
TLC taxi-zone identifiers for pickup and drop-off \\

\texttt{trip\_miles}, \texttt{trip\_time} &
Numeric &
Trip distance and trip duration \\

\texttt{base\_passenger\_fare} &
Numeric &
Base passenger fare recorded for the trip \\

\texttt{driver\_pay}, \texttt{tips}, \texttt{tolls} &
Numeric &
Financial fields associated with the trip \\

Shared/WAV flags &
Categorical &
Indicators relating to shared and wheelchair-accessible trips \\

\bottomrule
\end{tabular}
\end{table*}

% ============================================================
\begin{figure*}[t]
\centering

\begin{minipage}{0.48\textwidth}
\centering
\includegraphics[width=\linewidth]{hvfhv_size_by_year.png}
\caption{HVFHV dataset size by year.}
\label{fig:size}
\end{minipage}
\hfill
\begin{minipage}{0.48\textwidth}
\centering
\includegraphics[width=\linewidth]{distance_vs_fare_clean.png}
\caption{Trip distance and base fare.}
\label{fig:scatter}
\end{minipage}

\end{figure*}

% ============================================================
\begin{figure*}[t]
\centering

\begin{minipage}{0.48\textwidth}
\centering
\includegraphics[width=\linewidth]{trips_by_hour.png}
\caption{Trips by pickup hour.}
\label{fig:hour}
\end{minipage}
\hfill
\begin{minipage}{0.48\textwidth}
\centering
\includegraphics[width=\linewidth]{trips_by_day.png}
\caption{Trips by day of week.}
\label{fig:day}
\end{minipage}

\end{figure*}

% ============================================================
\begin{figure*}[t]
\centering

\begin{minipage}{0.48\textwidth}
\centering
\includegraphics[width=\linewidth]{top_pickup_zones.png}
\caption{Top pickup zones.}
\label{fig:zones}
\end{minipage}
\hfill
\begin{minipage}{0.48\textwidth}
\centering
\includegraphics[width=\linewidth]{pickups_by_borough.png}
\caption{Pickups by borough.}
\label{fig:borough}
\end{minipage}

\end{figure*}

\end{document}